\documentclass{article}
\usepackage{amsmath,amssymb}
\usepackage{xurl}
\usepackage[colorlinks=true,linkcolor=blue,citecolor=blue,urlcolor=blue]{hyperref}
% Shared commands for notes in docs/paper-gaps/.

\DeclareUnicodeCharacter{2080}{\ifmmode{}_{0}\else\textsubscript{0}\fi}
\DeclareUnicodeCharacter{2081}{\ifmmode{}_{1}\else\textsubscript{1}\fi}
\DeclareUnicodeCharacter{2082}{\ifmmode{}_{2}\else\textsubscript{2}\fi}
\DeclareUnicodeCharacter{2099}{\ifmmode{}_{n}\else\textsubscript{n}\fi}

\newcommand{\Lean}{\textsc{Lean}}
\ExplSyntaxOn
\NewDocumentCommand{\leanid}{m}{
  \ifmmode
    \textsf{\detokenize{#1}}
  \else
    \group_begin:
    \urlstyle{sf}
    \tl_set:Nn \l_tmpa_tl {#1}
    \tl_replace_all:Nnn \l_tmpa_tl {\_} {_}
    \exp_args:NV \path \l_tmpa_tl
    \group_end:
  \fi
}
\ExplSyntaxOff
\providecommand{\tr}{\operatorname{tr}}
\newcommand{\tnleanIssueBase}{https://github.com/LionSR/TNLean/issues/}
\newcommand{\tnleanPrBase}{https://github.com/LionSR/TNLean/pull/}
\newcommand{\tnleanDocBase}{https://sirui-lu.com/TNLean/docs/}
\newcommand{\ghissue}[1]{\href{\tnleanIssueBase#1}{GitHub issue \##1}}
\newcommand{\ghpr}[1]{\href{\tnleanPrBase#1}{PR \##1}}
\newcommand{\doclink}[1]{\href{\tnleanDocBase#1.html}{\path{#1}}}

% Dirac notation and the identity operator, as in blueprint/src/macros/common.tex.
% \providecommand keeps note-local definitions authoritative.
\usepackage{dsfont}
\providecommand{\ket}[1]{|#1\rangle}
\providecommand{\bra}[1]{\langle #1|}
\providecommand{\braket}[2]{\langle #1 | #2 \rangle}
\providecommand{\ketbra}[2]{|#1\rangle\!\langle #2|}
\providecommand{\Id}{\mathds{1}}
\providecommand{\one}{\mathds{1}}
\providecommand{\C}{\mathbb{C}}
\providecommand{\MN}[1]{M_{#1}(\C)}
\providecommand{\MD}{M_D(\C)}

% Verdict marker: \gapnote{<kind>}{<status>}. Kind is the policy
% classification (clarification, local-correction, scope-restriction,
% unfaithful, false-source, open-gap); status is open, wip (elimination
% underway), resolved, or historical. Machine-read by the site index;
% prints a verdict line.
\newcommand{\gapnote}[2]{\par\noindent\textbf{Verdict:} #1 (#2).\par}

\providecommand{\leanid}[1]{\textsf{\detokenize{#1}}}
\providecommand{\gapnote}[2]{\par\noindent\textbf{Verdict:} #1 (#2).\par}
\title{The Constant of $[H,H_L]$ in Hastings' Lemma~1}
\author{}
\date{October 22, 2026}
\begin{document}
\maketitle
\gapnote{false-source}{resolved}

\section{The source statement}

In the appendix of~\cite{Hastings2007AreaLaw} (arXiv:0705.2024v4,
\path{law16.tex}, line~794), the Hamiltonian $H=\sum_iH_{i,i+1}$, with
$\Vert H_{i,i+1}\Vert\le J$, is split at a cut as $H=H_L+H_B+H_R$, and the
source states: ``Let $b_L=[H,H_L]$. Note that $\Vert b_L\Vert\le J^2$.''

\section{A counterexample}

Take three qubits, $H_{12}=X_2$ and $H_{23}=Z_2$ (Pauli matrices on the second
qubit), so $J=1$, and let $H_L=H_{12}$. Then
\begin{align}
    b_L=[H,H_L]=[Z_2,X_2]=2iY_2,\qquad \Vert b_L\Vert=2>J^2 .
    \label{eq:hastings07_commutator_witness}
\end{align}
In general $[H,H_L]$ is a sum of commutators $[H_{i,i+1},H_{k,k+1}]$ of terms
on either side of the edge of the left part, each of norm up to $2J^2$.

\section{The statement proved}

For terms of range $R$, only the terms split by the region of $H_L$ contribute
to $[H,H_L]$. There are at most $N_R=2\cdot2^{2R}$ of them, each meets at most
$(R+1)2^{2R+1}$ terms of $H_L$, and each commutator has norm at most $2J^2$.
Hence
\begin{align}
    \Vert[H,H_L]\Vert\le2N_R(R+1)2^{2R+1}J^2,
    \label{eq:hastings07_commutator_bound}
\end{align}
and the same for $H_R$; for the buffer, $[H,H_B]=-([H,H_L]+[H,H_R])$. The
commutators are supported within distance $2R$ of the edges of the buffer.%
\footnote{\leanid{TNLean.MPS.AreaLaw.ChainHamiltonian.norm_commutator_leftPart_le}
in \path{TNLean/MPS/AreaLaw/Hastings/PartCommutators.lean}.}

\section{Conclusion}

The witness (\ref{eq:hastings07_commutator_witness}) shows that the printed
bound $\Vert b_L\Vert\le J^2$ is false, and (\ref{eq:hastings07_commutator_bound})
replaces $J^2$ by a constant depending only on the range. The bound enters
the later estimates of the proof of Lemma~1 only through this constant, so the
exponential decay of the error and the form of the area law are unchanged.

\bibliographystyle{alpha}
\bibliography{references}
\end{document}
